% Generated from validated Article Draft Result. Do not hand-edit; revise the structured draft instead.

Diffusion Models Beat GANsは、diffusion modelが画像生成のquality benchmarkをめぐる主要な研究系統として扱われていたことを示す。本稿では論文タイトルをそのまま普遍的な勝敗宣言として採用しない。benchmark設定や著者評価の範囲を越えて『diffusionがGANを全面的に置き換えた』と結論するのではなく、2021年にdiffusion qualityが独立した競争軸として明瞭になったという境界までを記録する。\autocite{src-11a149a699a0}


GLIDEは、diffusionをtext-guided generationとeditingへ接続する年末の技術境界として位置付けられる。ここで重要なのは、画像品質を競う生成器という見方から、language conditionを介して利用者の意図を画像生成へ渡すinterfaceへ問題が広がった点である。これによりdiffusionは単独の画像modelではなく、textとimageを結ぶ生成systemの一部として読めるようになった。\autocite{src-50244427230b}


Latent Diffusion Modelsは、別の方向から計算効率の問題へ踏み込む。画像空間そのものだけで処理を考えるのではなく、latent spaceを利用してgenerationのcomputational efficiencyを設計対象にする系統である。GLIDEのtext guidanceとLDMのlatent-space efficiencyは同じ改善ではない。条件付けのinterfaceと計算表現の選択という二つのleverが、年末のdiffusion研究に並んでいたことが重要である。\autocite{src-a4561c635fe6}


三つの資料を年次trajectoryとして並べると、2021年のdiffusionは『品質の高い新しい画像生成法』という一点から、quality、text guidance/editing、latent-space efficiencyという複数層へ展開したと読める。この構造は、翌年のimage-generation ecosystemでgeneration model、conditioning、distribution、runtime efficiencyが結び付く前段階にある。ただし、2022年に成立した製品・配布形態を2021年の論文へ遡及的に含めることはしない。\autocite{src-11a149a699a0,src-50244427230b,src-a4561c635fe6}


したがって2022年へのhandoffは、後年の成功を予言したという意味ではない。2021年末までに、diffusionをtextで制御する方向とlatent representationで効率化する方向が一次資料上に現れており、翌年に画像生成がより広い利用surfaceへ展開するための技術的部品が既に分離して見えていた、という限定された連続性である。2021年の価値は、完成したecosystemではなく、その構成要素が揃い始めたところにある。\autocite{src-50244427230b,src-a4561c635fe6}


\begin{claimboundary}
quality、generation/editing、efficiencyに関する具体的な優位性は原著者の評価であり、独立再現済みの事実として扱わない。本稿が確定するのは、これら三つの研究系統が2021年のdiffusion trajectoryを構成したことまでである。\autocite{src-11a149a699a0,src-50244427230b,src-a4561c635fe6}
\end{claimboundary}
